% =============================================================
%  Cables — what changes against an overhead line
%  Indicative ranges; the example values are the 230 kV XLPE cable
%  used on the power-exchange slides.
% =============================================================
\subsection{Cables}
\label{sec:cable}

A cable carries the same four parameters and the same representation, but
the geometry inverts the shunt branch. Conductor and screen are separated by
a few centimetres of solid dielectric rather than by metres of air, so the
capacitance per unit length rises by more than an order of magnitude while
the inductance falls, Table~\ref{tab:cable-ranges}.

\begin{table}[!ht]
  \centering
  \begin{tabular}{@{}l l r r@{}}
    \toprule
    \textbf{Parameter} & & \multicolumn{1}{c}{\textbf{OHL}}
      & \multicolumn{1}{c}{\textbf{Cable}} \\
    \midrule
    $R'$ & $\Omega$/km  & 0.1--1    & 0.1--1 \\
    $L'$ & mH/km        & $\approx 1$ & 0.31--1.42 \\
    $G'$ & \textmu S/km & 0.05      & $\approx 1$ \\
    $C'$ & nF/km        & $\approx 10$ & 180--420 \\
    \bottomrule
  \end{tabular}
  \caption{Indicative per-unit-length parameters, overhead line against
           XLPE cable~[17]}
  \label{tab:cable-ranges}
\end{table}

Three consequences follow from these values. Together they make a cable a
different design problem from a shorter overhead line.

\paragraph{Low wave impedance.} With $C'$ up and $L'$ down,
$\underline{Z}_w$ falls from a few hundred ohms to under a hundred. For a
230\,kV XLPE cable with $L' = 1.35$\,mH/km and $C' = 190$\,nF/km,
$\underline{Z}_w = 84\,\Omega$ against $266\,\Omega$ for the overhead line
of Section~\ref{sec:reference-line}. The natural power rises in proportion,
but it is not a loading that is reached: the charging current binds
first.

\paragraph{High charging power.} The charging reactive power per
unit length is $Q_c' = U^2 \cdot \omega \cdot C'$. For the same example
cable $\omega \cdot C' = 71.6$\,\textmu S/km, giving 3.79\,MVAr/km and a
charging current of 9.5\,A/km at 230\,kV. That current flows whether or not
the cable is loaded, it is drawn from both ends, and it consumes conductor
ampacity that is then unavailable for load current. Beyond a critical length
the charging current occupies the whole ampacity. Shunt reactors at
the ends --- and, for long routes, at intermediate points --- are therefore a
condition of operation.

\paragraph{Short electrical length.} The phase constant
$\beta = \omega \cdot \sqrt{L' \cdot C'}$ rises with the capacitance, so a
cable reaches a given $\beta \cdot l$ in a fraction of the route length an
overhead line needs. The example cable has a wavelength of 1041\,km against
4936\,km for the reference overhead line, so one kilometre of this cable has
the electrical length of 4.7\,km of that line. The lumped $\pi$ of
Section~\ref{sec:ac-two-port} therefore loses accuracy at a few tens of
kilometres rather than a few hundred, and the distributed form
of~\eqref{eq:ac-abcd-exact} applies at correspondingly shorter routes.

Where intermediate compensation is impossible --- a submarine crossing, or a
long cable with no accessible landing point --- these three together rule
out AC. The comparison is then between an AC and a DC scheme rather than
between a cable and an overhead line.

\FloatBarrier
